%% ========================================================================
%% Chapter 3: Input/Output (I/O) Basics
%% ========================================================================

\chapter{Input/Output (I/O) Basics}
\label{ch:basic-io}

\chapterhero{violet}{Learn how data streams work in the shell through file descriptors, redirection, pipelines, and efficient input-output processing techniques.}{\faIcon{exchange-alt}\quad stdin/stdout}{\faProjectDiagram\quad pipeline}{\faIcon{sign-in-alt}\quad redirection}

Imagine working at a photocopy shop. There are three work lanes: an input lane for receiving documents from customers, an output lane for returning copied documents, and a special lane for error reports when the machine fails. Linux uses the same idea for every program: there is a dedicated channel for input, a normal output channel, and an error message channel. This is the basic concept of I/O in Linux.

Unix philosophy states that each program should do one thing well, and that small programs can be combined to solve complex tasks. This combination is possible because all programs share a uniform I/O interface. The output of one program can directly become the input of the next program, like a production line in a factory.

Understanding I/O is required for almost every Linux administration task. Saving logs, filtering output, and separating error messages from normal data all depend on the mechanisms studied in this chapter.

\textbf{What You'll Learn}

By the end of this chapter, you will be able to:

\begin{enumerate}
    \item explain the concept of file descriptors and identify stdin, stdout, and stderr;
    \item redirect output to files using the \texttt{>} and \texttt{>>} operators and separate stderr with \texttt{2>};
    \item build multi-stage pipelines using the \texttt{|} operator with combinations of \cmd{grep}, \cmd{sort}, \cmd{uniq}, and \cmd{wc};
    \item use \cmd{tee} to send output to the terminal and a file at the same time;
    \item apply here-documents to create files with multi-line content directly from the command line.
\end{enumerate}

\begin{notebox}
All labs in this chapter use the directory \dir{~/lab-os/chapter03-io/}.
Create this directory before starting the first lab:
\begin{lstlisting}[language=bash]
mkdir -p ~/lab-os/chapter03-io
cd ~/lab-os/chapter03-io
\end{lstlisting}
\end{notebox}

%% ------------------------------------------------------------------------
\section{File Descriptor}
\label{sec:file-descriptors}

Whenever a program opens a file, the kernel gives it an integer called a \textit{file descriptor} (FD) as a handle for read and write operations. This number is lightweight and efficient: the program only needs to mention the number, and the kernel knows which file is meant.

The first three file descriptors are available to every process as soon as it starts. FD~0 is \textit{stdin} (standard input): the channel for reading input, connected to the keyboard by default. FD~1 is \textit{stdout} (standard output): the channel for normal output, directed to the terminal screen by default. The third is slightly different: FD~2, or \textit{stderr}, is dedicated to error messages and can be redirected separately from stdout, even when both normally point to the screen.

The separation between stdout and stderr is not just convention. When you build a pipeline such as \texttt{command1 | command2}, only stdout from \texttt{command1} flows to stdin of \texttt{command2}. Error messages from \texttt{command1} do not flow through the pipeline, so they do not contaminate the data processed by \texttt{command2}.

Linux stores information about open file descriptors for each process in the \dir{/proc/PID/fd} directory. Here, the shell variable \texttt{\$\$} refers to the PID of the currently running shell process.

\subsection*{Lab 3.1: Observe File Descriptors of an Active Process}

The goal of this lab is to see the file descriptors owned by the shell process and understand where each one points.

\begin{enumerate}
    \item Enter the working directory:
\begin{lstlisting}[language=bash, caption={Entering the working directory}]
cd ~/lab-os/chapter03-io
\end{lstlisting}

    \item View the file descriptors owned by the current shell:
\begin{lstlisting}[language=bash, caption={File descriptors of the active shell process}]
ls -l /proc/$$/fd
\end{lstlisting}
    \textit{Observe:} The first three entries (0, 1, 2) usually point to \file{/dev/pts/0} or a similar name. That is the pseudo-terminal representing your terminal session. All three point to the same place because input, output, and error are all connected to the same screen.

    \item View the current shell PID:
\begin{lstlisting}[language=bash, caption={Current shell PID}]
echo "PID shell ini: $$"
\end{lstlisting}

    \item Open a file for reading and inspect the new file descriptor that appears:
\begin{lstlisting}[language=bash, caption={Observing a new file descriptor}]
exec 3< /etc/hostname
ls -l /proc/$$/fd
exec 3<&-
\end{lstlisting}
    \textit{Notice:} After running \texttt{exec 3<}, file descriptor 3 appears in the list and points to \file{/etc/hostname}. After \texttt{exec 3<\&-}, that file descriptor is closed and disappears from the list.

    \item Check file descriptors for another process, for example init:
\begin{lstlisting}[language=bash, caption={File descriptors of process PID 1}]
sudo ls -l /proc/1/fd | head -10
\end{lstlisting}
\end{enumerate}

\begin{challengebox}
Redirect only stderr to a file while stdout still appears in the terminal. Run \cmd{find /etc -name "*.conf"}, which will produce two kinds of output: a list of found files (stdout) and ``Permission denied'' messages (stderr). Separate them with: \texttt{find /etc -name "*.conf" 2> error-find.txt}. Check \file{error-find.txt} with \cmd{cat error-find.txt}. How many permission errors were recorded?
\end{challengebox}

%% ------------------------------------------------------------------------
\section{Output Redirection}
\label{sec:io-redirection}

Output redirection operators instruct the shell to connect stdout or stderr to a file instead of the terminal. This is the most basic way to save command results.

The \texttt{>} operator redirects stdout to a file. If the file already exists, its contents are overwritten completely. If it does not exist, a new file is created. The \texttt{>>} operator also redirects stdout to a file, but in append mode: old content is preserved and new output is added at the end. This difference is crucial. Using \texttt{>} accidentally on a log file can erase history that cannot be recovered.

For stderr, use \texttt{2>} for overwrite or \texttt{2>>} for append. The number 2 is the file descriptor number for stderr. To combine stderr into the same stream as stdout, use \texttt{2>\&1}: this means ``redirect FD 2 to the same place as FD 1''. Order matters: write \texttt{> file 2>\&1}, not the reverse.

The special file \file{/dev/null} works like a trash bin. Everything written there disappears immediately. Use it to discard unneeded output, such as irrelevant ``Permission denied'' messages when searching across the entire system.

\subsection*{Lab 3.2: Save Command Output to Files}

The goal of this lab is to practice using \texttt{>}, \texttt{>>}, \texttt{2>}, and \texttt{2>\&1} operators in realistic situations.

\begin{enumerate}
    \item Save the file list from the working directory to a file:
\begin{lstlisting}[language=bash, caption={Saving output to a new file}]
ls -lh ~/lab-os/ > lab-list.txt
cat lab-list.txt
\end{lstlisting}
    \textit{Observe:} Output does not appear in the terminal because it has been redirected to a file. The \file{lab-list.txt} file contains exactly the same result that would normally appear on the screen.

    \item Append system information to the same file without deleting the old content:
\begin{lstlisting}[language=bash, caption={Appending output to an existing file}]
echo "=== System Info ===" >> lab-list.txt
uname -a >> lab-list.txt
date >> lab-list.txt
cat lab-list.txt
\end{lstlisting}

    \item Create log files that separate stdout and stderr:
\begin{lstlisting}[language=bash, caption={Separating stdout and stderr into different files}]
find /etc -name "*.conf" > konfig-ditemukan.txt 2> konfig-error.txt
echo "Files found: $(wc -l < konfig-ditemukan.txt)"
echo "Error lines: $(wc -l < konfig-error.txt)"
\end{lstlisting}
    \textit{Observe:} Two separate files are created: one contains paths of files successfully found, and the other contains permission error messages.

    \item Discard all error messages and save only successful output:
\begin{lstlisting}[language=bash, caption={Discarding stderr to /dev/null}]
find /etc -name "*.conf" 2>/dev/null > konfig-bersih.txt
wc -l konfig-bersih.txt
\end{lstlisting}

    \item Combine stdout and stderr into one log file:
\begin{lstlisting}[language=bash, caption={Combining stdout and stderr into one file}]
find /var/log -name "*.log" > semua-output.txt 2>&1
head -20 semua-output.txt
\end{lstlisting}
    \textit{Notice:} The \file{semua-output.txt} file contains a mix of found file paths and error messages in the order they occurred.
\end{enumerate}

\begin{challengebox}
Separate stdout and stderr into two different files with one command. Run \cmd{ls /etc /missing-directory}, using one existing directory and one nonexistent directory. Use: \texttt{ls /etc /missing-directory > result.txt 2> error.txt}. Check both files. Then repeat by combining both streams with \texttt{\&>}: \texttt{ls /etc /missing-directory \&> all.txt}. Compare \file{all.txt} with the manual combination of \file{result.txt} and \file{error.txt}.
\end{challengebox}

%% ------------------------------------------------------------------------
\section{Pipe and Pipeline}
\label{sec:pipes}

A \term{pipe} is a mechanism that connects stdout from one command directly to stdin of the next command, without an intermediate file. The pipe operator is \texttt{|}. Processes in a pipeline run at the same time: the command on the left starts producing output, and the command on the right processes it immediately without waiting for the left command to finish completely.

The power of a \term{pipe} is composition. Each Unix command does one job well. \cmd{sort} sorts lines. \cmd{uniq} removes adjacent duplicate lines, and the \texttt{-c} flag adds counts. \cmd{wc} counts lines, words, or characters. \cmd{head} and \cmd{tail} take specific parts of output. Combinations of these simple commands can answer complex questions.

As you learned in Chapter~2 with \cmd{grep} and \cmd{awk}, pipelines become more useful when combined with text manipulation commands. For example, to find the ten IP addresses that appear most often in a web log: read all lines with \cmd{cat}, extract the IP column with \cmd{awk}, sort with \cmd{sort}, count occurrences with \cmd{uniq -c}, sort again numerically with \cmd{sort -rn}, and take the top ten with \cmd{head -10}. All of that can be done in one command line.

Long pipelines can be split across multiple lines using a backslash \texttt{\textbackslash} at the end of a line for readability.

\subsection*{Lab 3.3: Build Multi-Stage Pipelines}

The goal of this lab is to build a pipeline step by step to analyze system data.

\begin{enumerate}
    \item Start with a simple command, then add one stage at a time:
\begin{lstlisting}[language=bash, caption={Step-by-step pipeline: count users per shell}]
cat /etc/passwd
\end{lstlisting}
    \textit{Observe:} Each line is one user entry with colon-separated columns.

    \item Extract the last column, the shell used, with \cmd{awk}:
\begin{lstlisting}[language=bash, caption={Extracting the shell column from /etc/passwd}]
cat /etc/passwd | awk -F: '{print $NF}'
\end{lstlisting}

    \item Sort the result:
\begin{lstlisting}[language=bash, caption={Sorting the shell list}]
cat /etc/passwd | awk -F: '{print $NF}' | sort
\end{lstlisting}

    \item Count how often each shell appears and sort from highest count:
\begin{lstlisting}[language=bash, caption={Counting each shell occurrence}]
cat /etc/passwd | awk -F: '{print $NF}' | sort | uniq -c | sort -rn
\end{lstlisting}
    \textit{Notice:} The output shows the count on the left and the shell name on the right, sorted from the most frequently used.

    \item Analyze disk usage: display the five largest directories in \dir{/var}:
\begin{lstlisting}[language=bash, caption={Five largest directories in /var}]
sudo du -sh /var/* 2>/dev/null | sort -h | tail -5
\end{lstlisting}
    \textit{Observe:} The \texttt{-h} flag for \cmd{sort} handles units such as K, M, and G correctly. Without this flag, \texttt{10M} can be treated as smaller than \texttt{9G}.

    \item Count running processes per user:
\begin{lstlisting}[language=bash, caption={Process count per user}]
ps aux | tail -n +2 | awk '{print $1}' | sort | uniq -c | sort -rn
\end{lstlisting}
    \textit{Observe:} \texttt{tail -n +2} skips the \cmd{ps aux} header line. The result shows which users run the most processes.
\end{enumerate}

\begin{challengebox}
Build a pipeline that combines techniques from Chapter~2, grep and awk, with sort and uniq. Analyze \file{/var/log/auth.log}, or \file{/var/log/syslog} if auth.log is not available. Find lines containing \texttt{Failed} using \cmd{grep}, extract a keyword from a chosen column using \cmd{awk}, count occurrences with \cmd{sort | uniq -c}, then sort the result. If neither file exists, create your own practice file with several lines containing the word Failed and analyze it using the same pipeline.
\end{challengebox}

%% ------------------------------------------------------------------------
\section{The tee Command and Advanced Techniques}
\label{sec:tee-advanced}

When you redirect output to a file with \texttt{>}, the output no longer appears in the terminal. This is a problem when running a long command that you want to monitor live and save at the same time. The \cmd{tee} command solves this: it reads from stdin and writes to stdout and to a file simultaneously. Its name comes from the shape of the letter T, which represents a data stream splitting into two branches.

Basic usage is: \texttt{command | tee file.log}. For append mode so an old file is not overwritten, add the \texttt{-a} flag: \texttt{command | tee -a file.log}. You can also write to several files at once: \texttt{command | tee file1.txt file2.txt}.

A here-document (\texttt{<<EOF}) is a way to provide multi-line input to a command directly from the command line, without creating a separate file first. The shell reads every line between the opening marker (\texttt{<<EOF}) and the closing marker (\texttt{EOF}) as input. The marker can be any word, as long as the opening and closing markers match. If the opening marker is enclosed in single quotes (\texttt{'EOF'}), variables inside the block are not expanded.

Here-documents are very useful for creating configuration files in scripts, writing multi-line messages, or providing input to commands such as \cmd{cat} and \cmd{tee}.

\subsection*{Lab 3.4: Output to a File and Terminal at the Same Time}

The goal of this lab is to use \cmd{tee} in pipelines and create files with here-documents.

\begin{enumerate}
    \item Use \cmd{tee} to view and save output at the same time:
\begin{lstlisting}[language=bash, caption={tee: output to terminal and file}]
df -h | tee disk-report.txt
echo "==="
cat disk-report.txt
\end{lstlisting}
    \textit{Observe:} The \cmd{df -h} output appears in the terminal while the command runs. The \file{disk-report.txt} file stores exactly the same content.

    \item Use \cmd{tee} in the middle of a long pipeline to tap the data stream:
\begin{lstlisting}[language=bash, caption={tee in the middle of a pipeline}]
ps aux | tee semua-proses.txt | grep -i ssh
\end{lstlisting}
    \textit{Notice:} The \file{semua-proses.txt} file stores all processes, the full output of \cmd{ps aux}, but the terminal shows only lines containing ssh.

    \item Append log entries with timestamps using \cmd{tee -a}:
\begin{lstlisting}[language=bash, caption={Repeated logging with tee -a}]
echo "$(date '+%Y-%m-%d %H:%M:%S') - disk check" | tee -a aktivitas.log
df -h | grep -v tmpfs | tee -a aktivitas.log
echo "$(date '+%Y-%m-%d %H:%M:%S') - finished" | tee -a aktivitas.log
cat aktivitas.log
\end{lstlisting}

    \item Create a configuration file using a here-document:
\begin{lstlisting}[language=bash, caption={Creating a file with a here-document}]
cat << 'EOF' > aplikasi.conf
# Application configuration
HOST=localhost
PORT=8080
MODE=production
LOG_LEVEL=info
EOF
cat aplikasi.conf
\end{lstlisting}
    \textit{Observe:} Single quotes around \texttt{'EOF'} prevent variables such as \texttt{\$HOST} from being expanded. Try removing the quotes and compare the result.

    \item Use a here-document with \cmd{tee} to write a file with superuser privileges:
\begin{lstlisting}[language=bash, caption={Writing a file with sudo using tee}]
cat << 'EOF' | sudo tee /tmp/test-sudo.conf > /dev/null
# Test file
test=success
EOF
cat /tmp/test-sudo.conf
\end{lstlisting}
    \textit{Notice:} The \texttt{>} operator is processed by the shell before \cmd{sudo}, so it does not receive root privileges. Combining \cmd{tee} with \cmd{sudo} is the correct way to write files that require superuser privileges.

    \item Create a repeated dated log using a complete pipeline:
\begin{lstlisting}[language=bash, caption={Dated log with pipeline and tee}]
LOGFILE="report-$(date +%Y%m%d).log"
{
    echo "=== Report: $(date) ==="
    echo "=== Disk ==="
    df -h | grep -v tmpfs
    echo "=== Memory ==="
    free -h
    echo "=== Top Processes ==="
    ps aux | sort -k3 -rn | head -5
} | tee "$LOGFILE"
echo "Log saved at: $LOGFILE"
\end{lstlisting}
\end{enumerate}

\begin{challengebox}
Build a complete pipeline that combines techniques from Chapter~2 and this chapter. Goal: create a complete process report that is saved and displayed in the terminal at the same time. Use \cmd{ps aux} to collect all processes, filter with \cmd{grep} to find only processes owned by your user (\cmd{grep \$USER}), sort by CPU usage with \cmd{sort -k3 -rn}, take the top 10 with \cmd{head -10}, and save to \file{top-proses.log} using \cmd{tee -a} so the log is not overwritten if run multiple times. Add a timestamp at the beginning with \cmd{echo "\$(date)"} before the main pipeline, also redirected to the same file.
\end{challengebox}

%% ------------------------------------------------------------------------
\section{Summary}

In this chapter, you learned:

\begin{itemize}[noitemsep,topsep=2pt]
    \item \textbf{File descriptors}: every process receives three automatic FDs: stdin (0), stdout (1), and stderr (2). All three can be redirected separately.
    \item \textbf{The \texttt{>} operator}: redirects stdout to a file by overwriting old content. Use \texttt{>>} to append without deleting existing content.
    \item \textbf{Separating stderr}: the \texttt{2>} operator redirects only error messages to a separate file. \texttt{2>\&1} combines stderr into the stdout stream.
    \item \textbf{/dev/null}: a special file that discards all input. Useful for suppressing irrelevant error messages.
    \item \textbf{Pipeline}: the \texttt{|} operator connects stdout from one command to stdin of the next. Processes run together and data flows without intermediate files.
    \item \textbf{tee}: sends output to the terminal while also saving it to a file. The \texttt{-a} flag enables append mode. It can be placed anywhere in a pipeline.
    \item \textbf{Here-document}: a \texttt{<\,<EOF ... EOF} block provides multi-line input directly to a command without needing a temporary file.
\end{itemize}

%% ------------------------------------------------------------------------
\section{Lab Assignment}

\textbf{General Instructions}: Complete all exercises independently. Record important steps, save screenshots when needed, and summarize the results as personal documentation.

\subsubsection*{Lab Assignment 3.1 I/O Redirection and Separation}
Run \cmd{find /usr -name "*.so" -size +1M}, which will produce two kinds of output: found file paths (stdout) and ``Permission denied'' messages (stderr). Do three things: (a) save stdout to \file{library-besar.txt} and stderr to \file{akses-error.txt} with one command, (b) count the lines in each file using \cmd{wc -l}, (c) repeat the command but this time discard stderr to \file{/dev/null} and save only stdout. Include screenshots and explain the difference in results.


\subsubsection*{Lab Assignment 3.2 System Analysis Pipeline}
Build pipelines to answer two system questions: (a) which shell is most commonly used by accounts on the system? Use \file{/etc/passwd} as the data source, extract the last column with \cmd{awk -F: '\{print \$NF\}'}, then count with \cmd{sort | uniq -c | sort -rn}. (b) How many active processes does each user have? Use \cmd{ps aux}, skip the header line, extract usernames, and count occurrences. Save both results to separate files using the \texttt{>} operator.


\subsubsection*{Lab Assignment 3.3 Creating an Automatic System Log}
Create a simple log directly from the terminal using the techniques you learned. Use a \texttt{\{ ... \}} block containing several commands: timestamp with \cmd{date}, disk usage with \cmd{df -h}, memory usage with \cmd{free -h}, and five processes with the highest CPU usage using \cmd{ps aux | sort -k3 -rn | head -5}. Pipe the entire block to \cmd{tee -a system-report.log} so the output appears in the terminal and is saved. Run it twice and verify that the log file contains two consecutive entries instead of being overwritten.


%
